\section{Representation quality is relative to a decision}
\label{sec:C05:sufficiency}
A representation is useful to an agent when it preserves distinctions that matter for future decisions under the available interaction policy. This criterion differs from visual similarity, image classification accuracy, or reconstructing every pixel. Let $O_t$ be the current multimodal observation and $z_t=f_\phi(O_t,H_t)$ an encoded state. An image encoder may discard lighting variation while preserving object identity; a computer-use agent may additionally need exact text, focus state, disabled controls, and the difference between a preview and a committed operation. These requirements depend on the action space and objective, so no representation can be declared generally sufficient from a single transfer benchmark.

For a one-step decision with loss $\ell(a,Y)$, define the Bayes risk available from information $X$ as $R^*(X)=\inf_\pi\mathbb E[\ell(\pi(X),Y)]$. If $z=f(X)$ is deterministic, then $R^*(X)\leq R^*(z)$: every policy on $z$ can be implemented on $X$ by composition, while the reverse need not hold. Equality holds when some optimal decision is measurable with respect to $z$. This elementary proposition supplies the right direction of comparison. Compression cannot improve the ideal information-theoretic optimum, although it can improve finite-data learning, optimization, and computational feasibility. Those practical benefits should not be confused with creating information absent from the observation.

For sequential control, a sufficient state must support relevant prediction and reward decisions across future actions. A representation trained on passive video can preserve correlations that disappear when the agent intervenes. A representation trained on one action distribution can omit details needed by another. The complete state $S_t$ includes external records, tool state, and permissions; a visual embedding is only one component. Session continuity therefore requires synchronizing encoded observations with event timestamps and action outcomes. An excellent screenshot representation paired with the wrong browser tab or stale environment state is still an incorrect decision input.

Evaluation should separate representation learning from policy learning. A frozen-encoder probe asks what a restricted decoder can recover. Fine-tuning asks whether the representation is a useful initialization under additional supervision. Closed-loop evaluation asks whether the integrated controller succeeds while its actions alter the observation distribution. These are different experimental interventions. A linear probe can fail even when information is present in a nonlinear code, and a powerful policy can compensate for a weak visual encoder through repeated tool queries. Reporting the intervention prevents either result from being overstated as a definitive property of the latent state.

